% ECONOMICS_PROBLEM: {"id": "D83-67", "title": "Competitive dynamics with dispersed imperfect information", "jel_code": "D83", "source_id": "OP-0584"}
% FIRST_POSTED: 2026-10-11
% ATTEMPT_STATUS: solved
% ENTRY_KIND: research_attempt
\documentclass[11pt]{article}
\usepackage[a4paper,margin=24mm]{geometry}
\usepackage[T1]{fontenc}
\usepackage{lmodern,microtype,amsmath,amssymb,amsthm,mathtools,xurl,hyperref}
\hypersetup{colorlinks=true,linkcolor=blue,urlcolor=blue,citecolor=blue}
\setlength{\parskip}{3pt}
\newtheorem{theorem}{Theorem}
\newtheorem{lemma}{Lemma}
\newtheorem{proposition}{Proposition}
\newcommand{\E}{\mathbb E}
\newcommand{\ind}{\mathbf1}
\newcommand{\R}{\mathbb R}
\title{Private beliefs, saving and exact recursive aggregation}
\author{JiHo Bae}
\date{October 11, 2026}
\begin{document}
\maketitle

\paragraph{Authorship and provenance.}
Author: JiHo Bae. Prepared with GPT Codex assistance; the precise serving
revision was not exposed. The complete original argument is retained in
the \href{https://github.com/jbaelaw/private-beliefs-storage-proof/blob/e0a1a66ce965da56b5a15203de2c631e61c20131/paper/Proof.tex}{immutable TeX source} and
\href{https://github.com/jbaelaw/private-beliefs-storage-proof/blob/e0a1a66ce965da56b5a15203de2c631e61c20131/paper/Proof.pdf}{proof PDF}. This submission adds catalogue metadata,
provenance and a completion statement to that argument. Its preparation
included exact-target and primary-source checks and separate mathematical
reviews. This is a complete proof claim submitted for independent expert
review; community acceptance remains pending.

\begin{abstract}
We solve the specified-class problem in catalogue entry D83-67 for finite-horizon competitive storage economies with dispersed private information. A finite Bayesian filter and backward optimization give household policies. An explicit common-prior population construction yields exact aggregate transitions and goods clearing, without an uncountable independence assumption. We characterize sufficient population histograms and construct the unique coarsest family by backward refinement. With rational finite data, the construction is polynomial in the explicitly supplied state graph. A private-information example shows why mean beliefs and even current aggregate actions can conceal the state needed for future aggregation.
\end{abstract}

\section{The economy and the scope of the result}
Problem D83-67 asks for equilibrium, belief dynamics and sufficient recursive statistics in a \emph{specified class} of non-Gaussian information economies. We specify that class here. Mankiw and Reis discuss heterogeneous updating and its extension to consumption and investment \cite[\S\S7.2, 7.4--7.5]{MR}. The finite-state formulation and aggregation theorem are our answer to the catalogue's editorial question, rather than a theorem attributed to that survey.

Dates are $\tau,\ldots,T$, with $T<\infty$. The class includes an economy starting at every initial date $\tau$ and every compatible initial population law. This convention permits uniform statements about recursive statistics; it does not assert that every population law is reachable from one fixed initial law.

There is a nonatomic unit population, finitely many household types $\theta$, a finite inventory grid $A_\theta\subset[0,\bar a]$ containing zero, hidden idiosyncratic states $Z_\theta$, and private observations $O_\theta$. An observation includes realized income $Y_\theta(o)>0$. The public state $x\in X$ follows an exogenous finite Markov chain with matrix $P_X$. Conditional on a public transition $x\to x'$, the joint private transition is
\[
 K_{\theta,t}^{x,x'}(z',o'\mid z).
\]
The public transition is independent of private hidden states given current public information. Saving does not affect either kernel. Household information consists of its identity, initial private label, private observations and incomes, its own chosen inventories, public states and spot prices, together with any independent private randomization it uses. All primitives and initial population laws are known. Nature's population-rotation seeds below are unobserved.

There is one spot good, whose normalized price is $p_t=1\in\Delta_1$. Households directly store goods: inventory $a'$ uses $a'$ current goods and produces $Ra'$ next period, with $R>0$. There is no separate bond or asset issuer. A feasible current plan satisfies
\begin{equation}\label{budget}
 c+a'=y+Ra,\qquad a'\in A_\theta,\qquad
 \underline c\le c\le\bar c,
\end{equation}
where $0<\underline c<\bar c<\infty$. Zero next inventory is always feasible, and at $T$ it is required. Household utility is $\E\sum_{t=\tau}^T\beta^{t-\tau}u_\theta(c_t)$, with $0<\beta\le1$ and continuous increasing $u_\theta$ on this compact consumption interval. Production is the direct storage technology in household plans. The class has no informative relative prices or financial trading; private information nevertheless affects intertemporal consumption.

Each type has a specified finite posterior alphabet $B_\theta\subset\Delta(Z_\theta)$. Compatible microstates form finite sets
\[
 S_t^x\subset\{\sigma=(\theta,a,q,y):a\in A_\theta,\ q\in B_\theta\}.
\]
They are closed under \emph{every} feasible inventory choice and every positive-probability private observation and public successor, using the Bayesian update below. Compatibility includes information already revealed by current income. Any probability vector on $S_\tau^{x_\tau}$ is an admissible initial law: Nature draws a private label $\sigma$ and then a hidden state with conditional law $q$. This gives a common prior consistent with the observed label and income.

At ties, select the smallest optimal next inventory at every microstate. All aggregation claims use this fixed equilibrium selection. If every Bellman maximum is strict, the selection describes all competitive-equilibrium control outcomes. For the computational statement, the grids, kernels, posteriors, returns, discount and incomes are rational, and rational utility values at every feasible consumption are supplied. A continuous increasing interpolation is part of the utility specification. No evaluation of an unspecified real-valued function is hidden in the algorithm.

\section{Individual beliefs and global optimization}
For $q\in B_\theta$ put
\begin{align}
 \ell^{x,x'}_{\theta,t}(q,o')
 &=\sum_{z,z'}q(z)K^{x,x'}_{\theta,t}(z',o'\mid z),\label{likelihood}\\
 \Phi^{x,x'}_{\theta,t}(q,o')(z')
 &=\frac{\sum_zq(z)K^{x,x'}_{\theta,t}(z',o'\mid z)}
               {\ell^{x,x'}_{\theta,t}(q,o')}.\label{filter}
\end{align}
The second expression is used only when its denominator is positive. Actual income is included in $o'$. Observed inventories and consumption are functions of already known resources, actions and private randomness. Prices are identically one. Thus none of these observations supplies an omitted hidden-state signal.

Set $V_T(\theta,a,q,y,x)=u_\theta(y+Ra)$. For $t<T$ define
\begin{multline}\label{bellman}
 V_t(\theta,a,q,y,x)=\max_{a'\text{ feasible}}\bigg\{
 u_\theta(y+Ra-a')\\
 +\beta\sum_{x'}P_X(x,x')
 \sum_{o'}\ell^{x,x'}_{\theta,t}(q,o')\,
 V_{t+1}\bigl(\theta,a',\Phi^{x,x'}_{\theta,t}(q,o'),Y_\theta(o'),x'\bigr)
 \bigg\}.
\end{multline}
Let $s_t^x(\sigma)$ be the selected maximizing inventory and let
$c_t^x(\sigma)=y+Ra-s_t^x(\sigma)$.

\begin{lemma}\label{individual}
The state $(t,x,\theta,a,q,y)$ suffices for the complete household problem. The selected policy is globally optimal against every adapted, possibly randomized, history-dependent saving plan.
\end{lemma}
\begin{proof}
Given the full private history, $q$ is the conditional law of the current hidden state. The Markov kernels and Bayes' rule give \eqref{likelihood}--\eqref{filter}, including income observations. Future public states have the displayed independent Markov law. Current feasible resources are known. Consequently the conditional distribution relevant to the next decision depends on private history only through the stated variables. The finite nonempty maximization in \eqref{bellman} bounds the payoff of every feasible action, including any private randomization. Backward iteration of this conditional inequality bounds every adapted plan; the selected rule attains equality at every step. This proves global optimality, without replacing the problem by first-order conditions.
\end{proof}

\section{A measurable population equilibrium}
An individual kernel alone does not justify exact aggregation of a continuum. The following construction specifies the joint population law as part of the class. It permits correlation across households, which the original problem does not prohibit.

\begin{lemma}\label{population}
Every admissible initial law has a common-prior realization in which all population transitions have their exact conditional masses, pathwise, and each household has the prescribed private transition kernel.
\end{lemma}
\begin{proof}
Take the population to be $I=[0,1)$ with Lebesgue measure $\lambda$. Partition it into half-open intervals of lengths $\mu_\tau(\sigma)$ and rotate this partition by an unobserved uniform $U_0$ modulo one. Each fixed household's initial private label has law $\mu_\tau$, while its realized population masses are exactly $\mu_\tau$. Knowing one's identity does not reveal the random label. In each realized label cell, use another independent uniform rotation to partition its normalized mass coordinate into hidden states with proportions $q(z)$. Every fixed household then has posterior $q$ after observing its label, and each realized initial hidden-state mass is $\mu_\tau(\sigma)q(z)$.

Thereafter group households by initial label and full \emph{exogenous} hidden-state and observation history. For a positive-mass cell $C$, use
\[
 v_C(i)=\frac{\lambda(C\cap[0,i])}{\lambda(C)}.
\]
Under normalized population measure on $C$, $v_C$ is uniform on $[0,1]$. After observing the next public state $x'$, partition
$(v_C(i)+U_{t,C})\bmod1$ by the inverse cumulative probabilities of
$K_{\theta,t}^{x,x'}(\cdot,\cdot\mid z)$, where $z$ is the cell's current hidden state. Fresh uniforms are independent of the past and the entire public chain. They may be indexed in advance by all possible finite history labels; only finitely many are needed.

For each realization, rotation preserves uniform mass coordinates, so every child cell has exactly its prescribed fraction of the parent mass. For each fixed household, conditional on the full past and $x'$, the selected fresh uniform is independent and uniform, even though its cell index is random. Its next state and observation therefore have the required conditional law. Half-open bins fix endpoint assignments. Empty cells need no update, and conditional laws on null histories receive arbitrary versions. These conventions preserve population integrals and each fixed household's almost-sure statements.

All cells remain finite unions of intervals: their child sets are preimages of finitely many intervals under rotated, piecewise affine mass coordinates. Their endpoints are measurable functions of finitely many past uniforms. Thus individual variables are jointly measurable in identity and the underlying finite product probability space, as are all aggregate integrals. The public finite chain can be included on that same space.

The cell indexing uses no saving actions. An individual deviation leaves its exogenous signal and income process unchanged. For the selected equilibrium, inventories and posteriors are functions of initial labels and observed histories, hence constant on each finer exogenous cell. This gives exact aggregation for the selected policy without claiming exact frequencies inside arbitrary identity-dependent deviation subsets.
\end{proof}

The construction also defines the original full-history belief
$b_{it}=\mathcal L(H_t\mid\mathcal F_{it})$. Take $H_t$ to be the finite-dimensional history of public states and underlying innovations, which determines all relevant population histories. This is a standard Borel space. If $O_i^t$ denotes the finite observed history under the selected policy, then for a Borel history event $D$,
\[
 b_{it}(D)=
 \frac{\int\ind\{H_t(\omega)\in D\}\ind\{O_i^t(\omega)=h\}\,d\mathbb P(\omega)}
 {\int\ind\{O_i^t(\omega)=h\}\,d\mathbb P(\omega)}
 \quad\text{on }O_i^t=h
\]
whenever the denominator is positive. Parameterized integration gives joint measurability in identity; null histories receive arbitrary versions. The original individual state includes identity as an economically irrelevant tag. Then $q$ is the own-hidden-state marginal of $b_{it}$, and $\theta,a,y$ are observed components. The map $(a_{it},b_{it})\mapsto(\theta,a,q,y)$ is a pointwise projection, not an assertion that a full history belief equals $q$.

Define the stochastic microstate matrix, conditional on $x\to x'$, by
\begin{equation}\label{microtransition}
 Q_t^{x,x'}(\sigma,\sigma')=
 \sum_{o'}\ell^{x,x'}_{\theta,t}(q,o')\,
 \ind\!\left\{\sigma'=
  (\theta,s_t^x(\sigma),\Phi^{x,x'}_{\theta,t}(q,o'),Y_\theta(o'))\right\}.
\end{equation}
\begin{theorem}\label{equilibrium}
Every member of the declared class has a selected competitive equilibrium. If $\mu_t$ is the cross-sectional law of $(\theta,a,q,y)$, its recursive law is
\[
 x_{t+1}\sim P_X(x_t,\cdot),\qquad
 \mu_{t+1}=\mu_tQ_t^{x_t,x_{t+1}}.
\]
Prices, aggregate consumption and next inventory are respectively
\[
 1,\qquad C_t=\sum_\sigma\mu_t(\sigma)c_t^{x_t}(\sigma),\qquad
 A_{t+1}=\sum_\sigma\mu_t(\sigma)s_t^{x_t}(\sigma).
\]
The $q$-marginal gives the reduced current-hidden-state posterior distribution, while the joint law retains its association with inventories and income. Full-history beliefs remain defined as above.
\end{theorem}
\begin{proof}
Initial calibration is
$\lambda\{\sigma_t=\sigma,Z_t=z\}=\mu_t(\sigma)q(z)$.
Exact cell splitting and Bayes' rule preserve this identity. Summing the calibrated child masses gives \eqref{microtransition} and the stated population law. In particular $\mu_t$ is a publicly known function of the initial law and public history. Lemma~\ref{individual} verifies global household optimality, and Lemma~\ref{population} verifies all beliefs under one common prior.

Integrating \eqref{budget} gives, in every realization,
\[
 C_t+A_{t+1}=\sum_\sigma\mu_t(\sigma)y
                       +R\sum_\sigma\mu_t(\sigma)a.
\]
Current storage is a production input, and $Ra$ is the output of previously owned storage. This is goods clearing with production; there is no omitted security market or intermediary. Terminal inventory is zero. Normalized prices, globally optimal plans, feasibility and consistent beliefs therefore form a competitive equilibrium. The projection of full-history beliefs suffices because individual future payoffs depend on the finite filter, and future aggregate laws follow the displayed public recursion.

Aggregate current endowment is independent of saving and of the private rotation realizations. It is a known function of the initial law and public history. If the catalogue's notation $w_t(x_t)$ is used literally, include that finite public history in $x_t$; the finite horizon preserves a finite Markov state. The net storage output is the production term, rather than an additional endowment counted twice.
\end{proof}

\section{Sufficient histograms and their smallest partition}
Write $g_t^x(\sigma)=(c_t^x(\sigma),s_t^x(\sigma))$. For a partition $\Pi_t^x$ of $S_t^x$, let $H_{\Pi_t^x}(\mu)$ be its vector of population masses. Calendar date and public state are retained. Sufficiency means that $(t,x,H_{\Pi_t^x}(\mu))$ determines current prices and aggregate controls, and the \emph{joint} conditional law of
\[
 \bigl(x',H_{\Pi_{t+1}^{x'}}(\mu Q_t^{x,x'})\bigr)
\]
throughout the declared class. Including the public successor prevents permutations of equally probable public shocks from concealing a failure of recursion.

\begin{theorem}\label{histograms}
A histogram family is sufficient if and only if, at every $t,x$, the following conditions hold:

\smallskip\noindent
(i) The output $g_t^x$ is constant within every current block.

\smallskip\noindent
(ii) For each $x'$ with $P_X(x,x')>0$ and each next block $D\in\Pi_{t+1}^{x'}$, the probability
$Q_t^{x,x'}(\sigma,D)$ is constant within every current block. At $T$, condition (ii) is omitted.

\smallskip\noindent
There is a unique coarsest sufficient family. It is constructed backward: terminal blocks group equal outputs; earlier blocks group equal current outputs and equal transition probability vectors to all previously constructed next blocks. Every sufficient partition refines this family. It therefore has the fewest bins, and $k-1$ independent histogram coordinates when it has $k$ bins. With the supplied rational finite data, household policies and this family are computable in polynomial time in the explicitly supplied state graph and tables.
\end{theorem}
\begin{proof}
If (i)--(ii) hold, let $g(B)$ and $\bar Q_t^{x,x'}(B,D)$ be the common output and row probabilities of block $B$. For current masses $m(B)$,
\[
 (C_t,A_{t+1})=\sum_Bm(B)g(B),\qquad
 m'(D)=\sum_Bm(B)\bar Q_t^{x,x'}(B,D).
\]
Together with $P_X$ and price one, these formulas give the required joint law.

Conversely, choose two microstates $\sigma,\widetilde\sigma$ in the same current block. The suffix economies starting from $\delta_\sigma$ and $\delta_{\widetilde\sigma}$ are admissible and have identical current histograms. Their outputs must agree. Conditional on any positive-probability public successor, their next histograms are deterministic row-probability vectors. Equality of the joint law forces these vectors to agree in every coordinate. This proves necessity using genuinely admissible initial economies.

The backward construction gives equivalence relations on finite sets and satisfies (i)--(ii). Any sufficient terminal partition refines the terminal output partition. Inductively, refinement of a constructed next partition implies that equality of probabilities to the finer blocks gives equality after summing to each constructed block. Current output equality then forces refinement of the constructed current partition. Backward induction proves coarseness, uniqueness and minimality. One histogram mass is recovered from the other $k-1$ by normalization. The minimality is among date- and public-state-indexed partition histograms, not arbitrary real encodings.

All filter calculations, Bellman comparisons and block-probability comparisons use rational arithmetic on the explicit finite graph. There are polynomially many such operations. If $D$ is the product of the supplied denominators, denominators of the expectations and values divide $D^{O(T+1)}$; their bit lengths are polynomial in the explicit unrolled input. Normalization in \eqref{filter} also uses the supplied rational posteriors. Sorting or pairwise comparison of the finite output and probability vectors implements refinement. The claim does not bound the size of a posterior alphabet constructed from an unrestricted signal model: that alphabet may grow exponentially or fail to be finite.
\end{proof}

\section{Private saving and a sharp compression obstruction}
\begin{proposition}\label{example}
The declared class contains strict private-saving economies in which mean beliefs, mean inventory and current income do not suffice. With three possible initial posterior states and zero initial inventory, the coarsest exact initial histogram has three bins and two independent masses.
\end{proposition}
\begin{proof}
Take two dates, $R=\beta=1$, inventory grid $\{0,1\}$, initial income three and terminal income $Z\in\{1,5\}$ with prior
$\nu=(\delta_1+\delta_5)/2$. An initial private erasure signal either reveals $Z$ or provides no information. In longer economies, further erasure signals may have type- and public-state-dependent revelation probabilities. Until terminal income is observed, unrevealing signals leave the posterior unchanged and revealing signals make it a point mass. Thus
$B=\{\nu,\delta_1,\delta_5\}$ is closed under the filter, including terminal income revelation.

For this example set $S_0=\{(0,q,3):q\in B\}$ and
$S_1=\{(a,\delta_z,z):a\in\{0,1\},\ z\in\{1,5\}\}$, suppressing the single preference type. These sets are closed under every feasible choice.

Use the continuous, increasing, concave piecewise linear utility on $[1,6]$ specified by
\[
 (u(1),u(2),u(3),u(4),u(5),u(6))
                  =\tfrac1{10}(0,7,11,14,16,18).
\]
At initial inventory zero, the values of saving zero and one are respectively $u(3)+\E_q u(Z)$ and $u(2)+\E_q u(Z+1)$. Their values, in posterior order $\nu,\delta_1,\delta_5$, are
\[
 (1.90,1.10,2.70)\quad\hbox{and}\quad(1.95,1.40,2.50).
\]
Hence the selected inventories are $1,1,0$, all strictly. Terminal inventory is required to be zero. Thus no tie convention affects this example.

Compare the initial populations
\[
 \mu^A=\delta_{(0,\nu,3)},\qquad
 \mu^B=\tfrac12\delta_{(0,\delta_1,3)}+
        \tfrac12\delta_{(0,\delta_5,3)}.
\]
They have the same preference type, mean inventory, current income, price and mean posterior $\nu$. The initial random-label construction gives the same hidden-income prior $\nu$ for every fixed household in both economies; the difference is private initial information. Nevertheless aggregate saving is one in the first economy and one half in the second. Aggregate consumption is two and two and a half. The proposed mean statistics fail.

Finally, current outputs separate $\delta_5$ from $\nu$ and $\delta_1$. The latter two both save one, but their terminal consumption masses differ: from $\nu$, consumption is two or six with equal masses; from $\delta_1$, it is two for everyone. Their probabilities to terminal output blocks differ. The backward criterion separates all three states, proving the claimed minimum. Current action bins alone therefore do not recover the state needed for future aggregation.
\end{proof}

This establishes equilibrium, private belief dynamics, recursive aggregate variables and exact minimal histogram sufficiency throughout the declared class. It answers D83-67's specified-class target. Informative relative-price feedback, general financial trade and unrestricted posterior spaces require different models and are not claimed here.

\section*{Completion Estimate}
100\%
The proof resolves the catalogue's equilibrium, aggregate-dynamics,
endogenous-belief and recursive-statistic requirements for the declared
private-information storage class. This is the author's complete proof
claim for that target, pending independent expert review.

\begin{thebibliography}{9}
\bibitem{MR} N. Gregory Mankiw and Ricardo Reis.
\emph{Imperfect Information and Aggregate Supply}.
Handbook of Monetary Economics, vol.~3A, chapter~5 (2011), pp.~183--229,
Sections 7.2, 7.4 and 7.5.
\url{https://doi.org/10.1016/S0169-7218(11)03005-X}.
Author-hosted copy: \url{https://personal.lse.ac.uk/reisr/papers/10-HofME.pdf}.
\end{thebibliography}
\end{document}
